% ===== PRÉAMBULE CANONIQUE — à coller VERBATIM en tête de CHAQUE .tex =====
\documentclass[11pt,a4paper]{article}
\usepackage[utf8]{inputenc}\usepackage[T1]{fontenc}\usepackage{lmodern}
\usepackage[french]{babel}
\usepackage{amsmath,amssymb,mathtools}\usepackage{icomma}
\usepackage{siunitx}\sisetup{locale=FR}  % \qty{}{}, \si{}, \ang{}
\usepackage{xcolor}\usepackage{colortbl}
\usepackage{tikz}\usepackage{pgfplots}\pgfplotsset{compat=1.18}
\usepackage[siunitx,europeanresistors]{circuitikz} % circuits (élec.)
\usepackage[most]{tcolorbox}\usepackage{enumitem}\usepackage{array,booktabs}
\usepackage[a4paper,margin=2.2cm]{geometry}
\usetikzlibrary{arrows.meta,calc,angles,quotes,positioning,patterns,%
  backgrounds,shapes.geometric,decorations.pathmorphing,decorations.markings,intersections}
\definecolor{primary}{RGB}{222,49,99}\definecolor{primarydark}{RGB}{150,22,55}
\definecolor{defcol}{RGB}{0,114,178}\definecolor{methcol}{RGB}{52,92,148}
\definecolor{excol}{RGB}{0,135,90}\definecolor{attncol}{RGB}{200,40,55}
\tcbset{boxbase/.style={enhanced,breakable,boxrule=0.9pt,arc=2.5pt,
  left=3mm,right=3mm,top=2.3mm,bottom=2.3mm,fonttitle=\bfseries,coltitle=white}}
% --- boîtes TUTORIEL (noms EXACTS, ne pas renommer) ---
\newtcolorbox{cle}[1][]{boxbase,colback=defcol!6,colframe=defcol,
  title={\textsc{À retenir}\ifx\relax#1\relax\relax\else\textmd{ : \itshape #1}\fi}}
\newtcolorbox{methode}[1][]{boxbase,colback=methcol!7,colframe=methcol,sharp corners,
  title={\textsc{Méthode}\ifx\relax#1\relax\relax\else\textmd{ --- \itshape #1}\fi}}
\newtcolorbox{exemplebox}[1][]{boxbase,colback=excol!5,colframe=excol,
  title={\textsc{Exemple}\ifx\relax#1\relax\relax\else\textmd{ : \itshape #1}\fi}}
\newtcolorbox{piege}[1][]{boxbase,colback=attncol!6,colframe=attncol,
  title={\textsc{Piège}\ifx\relax#1\relax\relax\else\textmd{ : \itshape #1}\fi}}
\newtcolorbox{remarque}[1][]{boxbase,colback=black!4,colframe=black!45,
  title={\textsc{Remarque}\ifx\relax#1\relax\relax\else\textmd{ : \itshape #1}\fi}}
% --- boîtes EXERCICE / CORRIGÉ (noms EXACTS, auto-comptées) ---
\newtcolorbox[auto counter]{exo}[1][]{boxbase,colback=primary!4,colframe=primary,
  title={\textsc{Exercice}~\thetcbcounter\ifx\relax#1\relax\relax\else\textmd{ --- \itshape #1}\fi}}
\newtcolorbox[auto counter]{corrige}[1][]{boxbase,colback=excol!4,colframe=excol,
  title={\textsc{Corrigé}~\thetcbcounter\ifx\relax#1\relax\relax\else\textmd{ --- \itshape #1}\fi}}
\newcommand{\vv}[1]{\vec{#1}}\newcommand{\ds}{\displaystyle}
\newcommand{\R}{\mathbb{R}}\newcommand{\N}{\mathbb{N}}\newcommand{\Z}{\mathbb{Z}}
% (un \usepackage{fancyhdr}+en-tête est possible ; il est ignoré côté web)
% =========================================================================
\begin{document}

\begin{center}
{\large\bfseries\color{primarydark} Thème 4 — Niveau Difficile}\\[-2pt]
{\color{primary}\rule{5cm}{1pt}}
\end{center}

\begin{exo}[un point sur deux sources]
Deux sources synchrones $S_1$ et $S_2$ émettent des ondes de longueur d'onde
$\lambda=\SI{2}{\centi\metre}$. Un détecteur est placé en $M$, avec $d_1=S_1M=\SI{10}{\centi\metre}$
et $d_2=S_2M=\SI{7}{\centi\metre}$.
\begin{center}
\begin{tikzpicture}[scale=1.0]
  \fill[black](0,0.9)circle(2.2pt)node[left,font=\footnotesize]{$S_1$};
  \fill[black](0,-0.9)circle(2.2pt)node[left,font=\footnotesize]{$S_2$};
  \fill[primary](5.6,0.1)circle(2.4pt)node[right,font=\footnotesize]{$M$};
  \draw[defcol,line width=0.9pt](0,0.9)--(5.6,0.1)node[midway,above,font=\footnotesize,black]{$d_1$};
  \draw[excol,line width=0.9pt](0,-0.9)--(5.6,0.1)node[midway,below,font=\footnotesize,black]{$d_2$};
\end{tikzpicture}
\end{center}
\begin{enumerate}[label=\alph*)]
  \item Calcule la différence de marche $\delta$.
  \item L'interférence en $M$ est-elle constructive ou destructive ? Justifie.
  \item En gardant les mêmes distances, donne une longueur d'onde $\lambda'$ qui rendrait
        l'interférence \emph{constructive} en $M$.
\end{enumerate}
\end{exo}

\begin{exo}[combien de franges brillantes sur l'écran ?]
Un dispositif de Young ($\lambda=\SI{600}{\nano\metre}$, $a=\SI{0.5}{\milli\metre}$,
$D=\SI{2}{\metre}$) éclaire un écran \textbf{large de \SI{2}{\centi\metre}}, centré sur la
frange centrale.
\begin{enumerate}[label=\alph*)]
  \item Calcule l'interfrange $i$.
  \item Les franges brillantes sont en $x_n=n\,i$ (avec $n$ entier relatif). Combien de franges
        brillantes voit-on sur l'écran (bord à bord : $|x|\le\SI{1}{\centi\metre}$) ?
\end{enumerate}
\end{exo}

\begin{exo}[deux haut-parleurs et un silence]
Deux haut-parleurs alimentés par le \emph{même} signal émettent un son de fréquence
$f=\SI{1700}{\hertz}$ (célérité du son $v=\SI{340}{\metre\per\second}$). On déplace un micro
devant eux.
\begin{center}
\begin{tikzpicture}[scale=1.0]
  \draw[black,line width=1pt](-0.15,0.7)rectangle(0.15,1.3);
  \draw[black,line width=1pt](-0.15,-1.3)rectangle(0.15,-0.7);
  \node[left,font=\footnotesize] at (-0.2,1){$S_1$};
  \node[left,font=\footnotesize] at (-0.2,-1){$S_2$};
  \fill[primary](5.4,0.4)circle(2.4pt)node[right,font=\footnotesize]{micro};
  \draw[defcol,line width=0.8pt](0.15,1)--(5.4,0.4)node[midway,above,font=\footnotesize,black]{$d_1$};
  \draw[excol,line width=0.8pt](0.15,-1)--(5.4,0.4)node[midway,below,font=\footnotesize,black]{$d_2$};
\end{tikzpicture}
\end{center}
\begin{enumerate}[label=\alph*)]
  \item Calcule la longueur d'onde $\lambda$ du son.
  \item Le micro peut-il ne \emph{rien} détecter ? Si oui, donne la plus petite différence de
        marche $\delta$ qui provoque ce silence.
  \item Explique en une phrase pourquoi un tel silence est possible.
\end{enumerate}
\end{exo}

\begin{exo}[deux ondes en quadrature]
Deux ondes de même amplitude $X$ et de même longueur d'onde se superposent. La première est en
\textbf{avance d'un quart de longueur d'onde} sur la seconde.
\begin{enumerate}[label=\alph*)]
  \item Calcule le déphasage $\Delta\varphi$ entre les deux ondes.
  \item L'amplitude de l'onde résultante vaut-elle $0$, $2X$, ou une valeur intermédiaire ?
        Justifie à partir de $\Delta\varphi$.
  \item \emph{(Pour aller plus loin.)} L'amplitude résultante de deux ondes de même amplitude
        $X$ déphasées de $\Delta\varphi$ vaut $A=2X\cos(\Delta\varphi/2)$. Calcule $A$ ici.
\end{enumerate}
\end{exo}

\begin{exo}[rouge ou bleu : quelles franges plus larges ?]
On réalise l'expérience de Young avec la même géométrie ($a=\SI{0.4}{\milli\metre}$,
$D=\SI{2}{\metre}$), d'abord en lumière rouge ($\lambda_R=\SI{700}{\nano\metre}$), puis en
lumière bleue ($\lambda_B=\SI{450}{\nano\metre}$).
\begin{enumerate}[label=\alph*)]
  \item Sans calcul, laquelle des deux donne l'interfrange le plus grand ? Justifie avec
        $i=\lambda D/a$.
  \item Calcule $i_R$ et $i_B$.
  \item Donne le rapport $i_R/i_B$ et vérifie qu'il vaut $\lambda_R/\lambda_B$.
\end{enumerate}
\end{exo}

\end{document}
